\section{Trait comparison, base forms and integral lifting}
\label{sec:traits}

Fix a package \(Y\to B\) from Proposition~\ref{bounded:preparation}, a
field-valued parameter with field \(K\), and a nonzero divisorial
valuation \(u\) on \(K\).  The generic parameter lifts on the good
isomorphism open, and properness extends it to
\(\Spec\cO_u\to B\).  We use \(\delta\) for a source cone, keeping
\(\sigma\) for the reference form.  Its stratum is \(O_\delta\), its
span lattice is \(N_\delta\), and its dimension is \(j_\delta\).

The four labelled formulas organize the argument.
Formula~\eqref{trait:TC} describes the actual normalized main trait
chart and its integral lattice. Formula~\eqref{trait:TE} separates
discrepancy into the base term, the toroidal order, and a nonnegative
excess; minimizing it gives \eqref{trait:T}, which identifies the
geometric minimum with the barred reference discrepancy. The two estimates in
\eqref{trait:TV} control retraction and variation along a labelled
cell. We then apply these calculations to construct the bounded-index
base form and prove integral lifting, the first two assertions of
Theorem~\ref{thm:bounded-fibre}. For lifting we return to the bounded
parameter and axis extensions; the auxiliary roots used in proving
\eqref{trait:TV} do not enter its denominator bound.

\subsection{Bounded saturation and the normalized trait chart}

Near the centre on the smooth base take boundary equations
\(t_1,\ldots,t_a\).  The base cone is \(\R_{\ge0}^a\), with lattice
\(\Z^a\).  A source cone over this centre maps onto it by \(p_*\).
Its stratum map is smooth: in toroidal charts the map of orbit tori
has full target rank, and characteristic zero makes it smooth.

We first arrange, simultaneously on every mapped face span, that
these maps are surjective on lattices.  Take \(h\)-th roots of the
base axes, where the fixed integer \(h\) is divisible by the indices
of all the finitely many image lattices in their saturated image
spans, and normalize the main pullback.  The base lattice is replaced
by \(h\Z^a\), and each source lattice by its preimage under \(p_*\).
Since \(h\Z^a\) is contained in every relevant full image lattice,
the new map is onto, including for every face mapping onto the base
cone.  For faces over smaller base faces the same assertion holds in
their image span.  Substitution of powers of the new axes, followed
by monoid saturation, proves the chart description.  Rational cone
shapes and their order functions pull back unchanged.  So does the
relative logarithmic determinant, as is immediate in rational
character spaces.

There can be several new strata of the same cone shape.  The new
source is nevertheless strict and simplicial: locally its different
boundary germs lie over different old boundary germs and have the
same real cone configuration.  We use finitely many such charts
along the fixed base strata.  The Kummer bases are smooth, and the
models remain projective over them.  Every old cone option is
available over every lifted base point.  Indeed, the axis cover is
finite flat; the main pullback before normalization meets every
point of the fibre product, since the universal generic fibre is
geometrically integral.  The same conclusion follows directly from
the local saturated monoid calculation.  Normalization is finite
and surjective onto that main pullback.

On an actual parameter this operation requires at most a bounded
finite field extension.  We prolong valuations with their old
scale, as in \eqref{bounded:finite-extension}.  Hence neither the
proposed discrepancy minimum nor its pulled-back order formula
changes.  From now on \(Y,B,t_j\) refer to these new data, and all
ordinary-fibre labels are taken \emph{after} this normalization.
Set
\[
 b_j=u(t_j)>0,\qquad \tau=u(\pi)>0,\qquad
 k_0=k(u),\qquad \ell=b/\tau\in\Z_{>0}^a,
\]
where \(\pi\) is a uniformizer of \(\cO_u\).  When \(a=0\), the
axis and relation lists are empty; all the arguments below include
this case.

Complete the trait, choose a coefficient section fixing \(\C\), and
extend the residue field to an algebraic closure \(\kappa\).  Thus
\[
 R=\kappa[[\pi]],\qquad K_R=\kappa((\pi)).
\]
Let \(s\in B_\kappa\) be the resulting geometric closed point.  Over
\(R\) take the normalization of the main component of the pullback
of \(Y\), whose generic field is that of the geometrically integral
generic model after scalar extension.  The following calculation
identifies its actual local branches.

At a geometric closed point \(y\in O_\delta\) over \(s\), put
\(M=N_\delta^\vee\) and \(P=\delta^\vee\cap M\).  Boundary
coordinates pull back to \(X^{p_j}\) times units, with
\(p_j=p^*e_j\).  Surjectivity of the lattice map splits its dual
inclusion integrally:
\begin{equation}
 M=p^*\Z^a\oplus M_{\mathrm{rel}}.
 \label{trait:primitive-splitting}
\end{equation}
Given prescribed unit multipliers on the \(p_j\), assign value one
on a basis of \(M_{\mathrm{rel}}\) and extend multiplicatively.
This produces a unit-valued character on \(M\) using only integral
powers of units.  In particular it does not extract roots of their
residues, even at the generic point of a labelled stratum over a
nonclosed residue field.

The substitution \(X^m\mapsto U(m)X^m\) preserves all monoid
relations and is an automorphism of the complete local ring.  On
its cotangent monoid generators it is multiplication by nonzero
residue scalars; the smooth variables can be fixed.  Successive
approximation gives surjectivity, and a surjective endomorphism of
a Noetherian ring is injective.  This also works on a singular
toric chart, where monoid atoms are not independent coordinates.
The old smooth base parameters can simultaneously replace the
appropriate smooth source parameters, since their transverse
Jacobian has full rank.  The combined cotangent matrix is block
triangular and invertible.  Make these changes on the source
first; after pullback, absorb the remaining trait units using
\eqref{trait:primitive-splitting} again.  The latter multipliers
come from the coefficient ring of the trait.  Eliminating the
smooth base parameters now gives precisely
\[
 X^{p_j}=\pi^{\ell_j}\quad(1\le j\le a),
\]
with \(n-j_\delta+a\) free smooth parameters.  All these changes
preserve the face ideals.  If a generic strict toroidal chart is
given only by an etale chart, the calculation may first be made
after a separable residue extension; the preserved face ideals
and the uniqueness below then descend it.

The dual cone and lattice of the normalized main chart are
\begin{equation}
 \begin{split}
 D_\delta&=\{(x,e)\in\delta\times\R_{\ge0}:p_*x=\ell e\},\\
 \Lambda_\delta&=\{(x,e)\in N_\delta\oplus\Z:p_*x=\ell e\}.
 \end{split}
 \tag{TC}\label{trait:TC}
\end{equation}
To verify normalization and completion in this assertion, consider
the character-group pushout
\[
 G=(M\oplus\Z)/
 \langle(p^*c,-\ell\cdot c):c\in\Z^a\rangle.
\]
By \eqref{trait:primitive-splitting}, it is
\(M_{\mathrm{rel}}\oplus\Z\), hence torsion-free.  Let \(S\) be
the image of \(P\oplus\N\) in \(G\), and \(\overline S\) its
saturation.  A point \(x_0\in\relint\delta\) with
\(p_*x_0=\ell\) exists because an onto cone map takes relative
interior onto relative interior.  The functional \((x_0,1)\) is
strictly positive on every nonzero image in \(S\).  Thus
\(\overline S\) is sharp.  The relation subspace meets the product
interior, so \(D_\delta\) spans
\(\Lambda_\delta\otimes\R\); no further quotient lattice is needed.

Localizing the binomial algebra at all monomials gives the group
algebra of \(G\), which has one torus component.  Its closure is
\(\kappa[S]\), and its finite normalization is
\(\kappa[\overline S]\).  Other components of the unsaturated
binomial equations, if any, lie on the monomial boundary.  Sharpness
gives one monomial origin.  The completion of the normal monoid
algebra there, with the free smooth parameters adjoined, is a
normal domain: the algebra is excellent, and a positive integral
grading gives the domain property for its monoid series.  Hence
there is one normalized main formal branch at this origin.

This branch is the one in the algebraic trait model.  Its generic
fibre lies in the good open, is geometrically integral, and is not
contained in any boundary component.  Boundary monomials are
therefore nonzerodivisors on the main model, and flat completion
preserves that fact.  No completed main branch is supported in
the monomial boundary.  The preceding torus closure calculation
then lists exactly the main completed branches; excellence
identifies their finite normalizations with completion of the
algebraic normalization.  This proves \eqref{trait:TC} for the
actual model, rather than merely for an abstract formal binomial
quotient.

\subsection{Monomial valuations, labels and algebraic ray divisors}

Faces of \(D_\delta\) meeting \(e>0\) correspond exactly to faces
\(\beta\le\delta\) mapping onto the base cone, and their relative
interiors correspond.  In the slice \(e=\tau\), a rational point
\(x\) defines a monomial valuation \(w_x\): take the minimum of
the weights \((x,\tau)\) in a monoid series, giving the free smooth
parameters weight zero, and then restrict to the function field.
This is a valuation also on a face.  The support of a nonzero
series generates a monomial ideal in a finitely generated monoid,
so finitely many exponents from that support give its minimum
against every nonnegative weight.  The zero-weight monomials and
smooth parameters form the coefficient part; products of two
nonzero initial forms remain nonzero in its series domain.  The
same argument proves that evaluation of every rational function
is continuous on the entire closed slice: write it as a quotient
of regular series and subtract the two finite minima.  It applies
to real as well as rational weights.  Moreover,
\begin{equation}
 (mx,m\tau)\in\Lambda_\delta
 \quad\Longrightarrow\quad mw_x\text{ is integral-valued}.
 \label{trait:integral-weight}
\end{equation}

If \(x\in\relint\beta\), label this valuation by the incident
irreducible component of \(O_\beta\cap Y_s\), taken geometrically
in the \emph{ordinary} parameter fibre.  At a deeper point in its
closure, the specified face determines its local branch, by
strictness.  The ordinary orbit-closure fibre germ for that face
is irreducible: base boundary equations vanish on the orbit
closure, and the smooth base parameters can be eliminated.  It
meets its open stratum densely.

The corresponding face centre in \eqref{trait:TC} has exactly this
ordinary-fibre branch.  The relevant comparison is integral:
\begin{equation}
 \Lambda_\delta/\Lambda_\beta\simeq N_\delta/N_\beta,
 \qquad
 N_\beta=N_\delta\cap\operatorname{span}_{\R}\beta.
 \label{trait:integral-quotient}
\end{equation}
Indeed, the map \((n,e)\mapsto n\bmod N_\beta\) has kernel
\(\Lambda_\beta\).  Given \(n\in N_\delta\), the already arranged
surjectivity \(p_*N_\beta=\Z^a\) supplies \(b'\in N_\beta\)
with \(p_*b'=-p_*n\).  Then \((n+b',0)\) lifts its class.  The
quotient cones agree as well.  For \(n\in\delta\), choose \(e\)
large enough that \(\ell e-p_*n\) is in the base cone, and choose
\(b'\in\beta\) mapping to this difference.  The point
\((n+b',e)\in D_\delta\) projects to the class of \(n\), and
the reverse containment is immediate.  For an empty axis list
the same proof simply uses the product cone.

Thus the quotient orbit-closure monoid has the original integral
lattice, cone and free smooth coordinates.  There is no finite
orbit cover at the generic ordinary-fibre component.  The old
unit-character changes preserve face ideals and are formal
automorphisms of these quotients; the subsequent trait
multipliers reduce to constants on the face centre, where the
axes and \(\pi\) vanish.  Hence the centre maps dominantly to
the stated ordinary branch.  At an interior point, where
\(\delta=\beta\), its free smooth coordinates are precisely
the coordinates of that branch.  Formula
\eqref{trait:integral-quotient} proves compatibility at every
deeper incident point.

We next show that rational \(w_x\) is an intrinsic algebraic ray
order with this label.  Choose effective Cartier multiples of
the boundary primes.  These are available on the strict
simplicial source and their characters span the dual of the
cone over \(\Q\).  For every positive boundary weight, compare
an appropriate positive power of its equation with a power of
\(\pi\) having the same weight at \((x,\tau)\).  The corresponding
two-generated ideals impose the equal-weight walls.  A zero
boundary weight already gives a face wall.  Together the walls
isolate the ray through \((x,\tau)\), because the boundary
equations give separate axis multiples on the simplicial cone.
These ideals are algebraic: use the Cartier divisor ideals
and \((\pi^q)\); changing a local generator by a unit does not
change either the ideal or its wall.  The choices can be used
consistently wherever this face and ray occur.

Their normalized blowups are projective vertical modifications,
and do not change the generic fibre.  In the completed chart,
adjoining the blowup ratios and saturating is the toric
subdivision by those walls.  It extracts the desired ray prime.
That prime remains irreducible after completion of the original
chart.  To see the last point explicitly, in the algebraic
toric modification its image is the orbit closure of the
supporting old face.  The generic fibre over this closure is
geometrically integral: the orbit-lattice map is the projection
modulo the larger face span, which is onto with torsion-free
kernel.  The supporting orbit closure is still integral on
completion.  Flat base change makes any newly appearing minimal
prime contract to the old ray prime.  It therefore meets the
inverse image of an open where this map of orbit closures is
flat, and on that open dominates the completed supporting
closure.  Geometric integrality of the generic fibre permits
only one such component.

The minimum-weight rule for monoid series is the scaled order
of this prime.  One precise verification is to make a smooth
toric refinement retaining the ray.  Powers of the ideal of
its ray divisor push down to the monomial ideals defined by
the corresponding order cutoffs.  Flat completion preserves
these pushforwards, and membership in the completed ideals is
termwise membership.  The ray divisor remains a single reduced
Cartier divisor there; the excellent completion is a regular
map.  This proves the order calculation on all series and
therefore on rational functions.

Contract the completed ray prime to the algebraic modification
over \(R\).  Its height is one: it contains the nonzero vertical
parameter, whereas flatness bounds the height of its
contraction by the height-one prime above it.  Its image
dominates the labelled ordinary-fibre component, by the
orbit calculation.  It is the unique algebraic ray prime
with the given label and weights.  Indeed, at an interior
point of that component the normalized main chart has one
branch and the same free stratum coordinates.  On its ray
modification there is one ray prime dominating that local
component.  Properness brings every global prime candidate
dominating the component above this chosen point.  Faithful
flatness lifts its generic point to the completed modification;
the equal-cut unit conditions and positive vertical order
place it in that single ray closure.  Its height-one
contraction forces all candidates to coincide.  Equivalently,
one can work at the generic component, using the primitive
unit splitting and \eqref{trait:integral-quotient} to exclude
both unit-root branches and orbit covers.

The orders computed at different interior points, or at
incident deeper points with this same label, consequently
agree, with the scale determined by \(w_x(\pi)=\tau\).
This is uniqueness of the toric ray order; a valuation with
additional positive excess can of course have the same
boundary weights.  Paths between labelled cells will use
only comparable strata with incident ordinary-fibre
components.  A segment, including its incident endpoints,
can be evaluated in the one full chart at the deeper labelled
point.  All labels thus come from the fixed parameter families.

These identifications also commute with a root of the
uniformizer.  If \(\pi=\rho^d\), then
\begin{equation}
 \Lambda_d=\{(n,e):p_*n=d\ell e\},\qquad
 \Lambda_d\longrightarrow\Lambda_\delta,\quad
 (n,e)\longmapsto(n,de).
 \label{trait:root-lattice}
\end{equation}
The weight \((x,\tau/d)\) maps to \((x,\tau)\), so the labelled
ray order restricts with its original scale.  The quotient
calculation \eqref{trait:integral-quotient} is independent of
\(\ell\), proving preservation of the ordinary-fibre label.

\begin{remark}
The preliminary saturation and the subsequent choice of
labels cannot be omitted.  On the open set
\(y(y^2+x)\ne0\), the map
\[
 t=x^2(y^2+x)
\]
has boundary lattice map multiplication by two at the
ordinary component \(x=0\).  Its generic fibre is
geometrically integral: with \(z=xy\), its equation is
\(z^2=t-x^3\).  After \(t=\pi^2\), normalization introduces
\(q=\pi/x\), with
\[
 q^2=y^2+x,\qquad \pi=xq.
\]
The special fibre over that old component has two primes,
\(q=y\) and \(q=-y\).  Both have boundary weights
\(w(x)=w(\pi)=1\), \(w(y)=0\), but their orders differ on
\(q-y\).  The bounded Kummer normalization separates them
into different new strata, each with primitive boundary
map.  Taking the labels on that new model is exactly what
makes the uniqueness proved above applicable.
\end{remark}

\subsection{Logarithmic coordinates and the excess formula}

Near a chosen stratum point, take equations for effective
Cartier multiples of the boundary primes.  Select
\(j_\delta-a\) of their characters which rationally complement
the base axes, and add regular units with independent
relative stratum differentials.  This gives a list of
\(n\) functions \(z=(z_1,\ldots,z_n)\).  In the complete
toroidal chart the character block and the stratum-coordinate
block are invertible, so
\(d\log z_1\wedge\cdots\wedge d\log z_n\) is a relative
logarithmic determinant frame.  Finitely many Zariski
neighbourhoods along the finitely many strata suffice.
Their functions have uniformly bounded rational
presentations in the fixed family embeddings: use finitely
many affine neighbourhoods and express their regular
functions by bounded-degree ratios whose denominators are
nonzero there.  A chart used by the main trait closure has
denominators surviving generically.  Hence these are actual
rational functions on the corresponding generic model.

The ratio
\begin{equation}
 F=\frac{\eta}
 {(d\log z_1\wedge\cdots\wedge d\log z_n)^{\otimes r_0}}
 \label{trait:frame-ratio}
\end{equation}
is toroidal locally, of order function \(r_0\phi\).  Algebraically,
a suitable power becomes a unit after dividing by the
corresponding boundary equations.  The retained fibrewise
differential identities and nonzero denominators ensure
that this assertion specializes to the prescribed form in
\(L\), including for nondominant parameter maps.

The labelled valuation \(w_x\) is the Gauss valuation on
\(K(z)\), with \(z\)-values linear in \(x\) and the smooth
units of weight zero.  First suppose \(x\) is rational and
interior to the cone.  Modulo the relations
\(p_j=\ell_j[\pi]\), the complementary boundary characters
and \([\pi]\) are rationally independent.  Their angular
residues, obtained by clearing denominators with powers,
therefore have \(j_\delta-a\) independent character
directions.  Their unit factors reduce to units in the
stratum parameters.  The remaining \(z_i\) reduce to
functions of the free smooth parameters with independent
differentials.  These two blocks give \(n\) algebraically
independent angular residues over the base residue field.
One can take a root of \(\pi\) to make each rescaled \(z_i\)
have weight zero, or take powers of the angular expressions
to stay over the original field.  This proves Gauss without
cancellation.  Continuity of series evaluation extends the
Gauss formula to faces and real weights in the fixed chart.

In particular, for rational \(x\), the restriction of
\(w_x\) to \(L\) is divisorial.  Its value group is discrete,
its residue field has the \(n\) independent additional
residues just exhibited over the divisorial base residue
field, and \(L/K(z)\) is finite.  The rank-one Abhyankar
criterion gives divisoriality.  The same argument applies
at a face, using its chart, and when some rounded
coefficients in the final construction below vanish.

Conversely, a rational divisorial \(w\) on \(L\) over \(u\)
has a centre with rational boundary weights \(x\) in one of
these slices.  It can be followed on the geometric trait
model with its scale unchanged.  Indeed,
\(\widehat L_w\) contains \(\widehat K_u\); a coefficient
field for the former can be chosen extending the chosen
one for the latter by lifting a residue transcendence
basis and then using separable Hensel lifting.  In a
Laurent-series presentation, extend this coefficient
field to an algebraic closure and thereby embed
\(\kappa((\pi))\) compatibly.  The scalar-extension function
field embeds as well: regularity of \(L/K\), and the
\(n\) independent additional residues from relative
Abhyankar equality, ensure that transcendence degree is
not lost after the algebraic base-residue extension.
The prolongation has the same value scale.  Its centre
selects a component label; use a closed specialization
within that component when choosing a chart, and let
\(w_x\) be the resulting retraction.  These large residue
and completion fields are used to calculate charts and
values of functions, never to define absolute form
discrepancies.

\begin{proposition}[Trait comparison]
\label{prop:trait-comparison}
For every parameter map allowed above and every nonzero
divisorial base valuation \(u\), the barred reference
discrepancy is
\begin{equation}
 \bar A_\sigma(u)=A_\zeta(u)+\mu(b).
 \tag{T}\label{trait:T}
\end{equation}
For a divisorial \(w\) over \(u\), with labelled monomial
retraction \(w_x\), there is an excess satisfying
\begin{equation}
 A_\sigma(w)=A_\zeta(u)+\phi(x)+E(w),\qquad
 E(w)\ge0,\qquad E(w_x)=0.
 \tag{TE}\label{trait:TE}
\end{equation}
Let \(\omega\) be any other form with effective boundary
on \(G\), and compare it to \(\sigma\) at a common tensor
index \(p\), writing
\[
 A_\omega(v)=A_\sigma(v)+\frac{v(H)}p.
\]
There is a uniform constant \(C_2\) such that
\begin{equation}
 \begin{split}
 |w(H)-w_x(H)|&\le C_2pE(w),\\
 |w_x(H)-w_{x'}(H)|&\le C_2p\|x-x'\|.
 \end{split}
 \tag{TV}\label{trait:TV}
\end{equation}
The second inequality holds along every segment in one
labelled slice cell, including incident endpoints, in
fixed cone coordinates.  The rational monomial valuations
in these assertions are the compatible algebraic labelled
ray orders constructed above.  Formulae \eqref{trait:T}
and \eqref{trait:TE} do not require absolute log
canonicity of \(\sigma\).
\end{proposition}

\begin{proof}
First compute the excess using finitely generated fields.
After a finite base extension by a root uniformizer
\(\widetilde\pi\), clear the ratios
\(w(z_i)/u(\pi)\) and put
\(z'_i=\widetilde\pi^{-h_i}z_i\), of value zero.
Prolong all valuations with the old scale.  Form
discrepancies are unchanged under this finite extension.
On the smaller rational field in the variables \(z'\)
over the extended base field, extract the prolonged base
prime and take, near its generic point, a smooth base
model times the smooth unit torus with coordinates
\(z'_i\).  Denote this model by \(T\), its reduced base
prime pullback by \(P_0\), and the prolonged restriction
of \(w\) by \(v\).  Then
\begin{equation}
 E(w)=A_{T,P_0}(v).
 \label{trait:excess-torus}
\end{equation}
To check the equality, use the finite-extension discrepancy
formula for \(L/K(z)\), the ratio
\eqref{trait:frame-ratio}, and invariance of the relative
\(d\log z\) determinant under rescaling by base
functions.  They reduce the absolute form to
\(\zeta\wedge d\log z'_1\wedge\cdots\wedge d\log z'_n\).
At the base prime, the divisor of \(\zeta\) contributes
exactly \(A_\zeta(u)\) after the discrepancy of the
reduced-fibre pair is separated off.  The contribution of
\(F\) is \(\phi(x)\).  This calculation uses the
homogeneous scale throughout and proves
\eqref{trait:TE}.  The smooth reduced pair \((T,P_0)\)
is lc, so \(E(w)\ge0\); for \(w_x\), the established
Gauss rule is the order of \(P_0\) with its scale,
giving \(E(w_x)=0\).  Every cone option occurs, and the
minimum defining \(\mu(b)\) has a rational minimizer
with a label.  Minimizing \eqref{trait:TE} therefore
proves \eqref{trait:T}, including for the universal
generic family.

We now prove \eqref{trait:TV}.  The combined degree of
the generic zero and pole divisors of \(H\), with
multiplicities, is at most \(C_1p\).  This remains a
uniform bound after the preliminary bounded Kummer
preparation.  Its normalized models are finite over
the old base pullbacks; over the good opens they give
the prescribed scalar-extended generic models mapping
to \(G\).  Pulling back the bounded polynomial cuts
on \(G\) bounds degrees in fixed projective embeddings
of these finitely many families.  Intersection degrees
of fixed line bundles in finite type are bounded after
stratifying by flatness.  The generic fibres used are
geometrically normal, both by preparation and, off the
base boundary, by eliminating smooth parameters in
the toroidal charts.

For the first inequality take the root uniformizer
large enough that
\(x/u(\widetilde\pi)\in N_\delta\).  This auxiliary
root degree is allowed to be unbounded.  If
\(\pi=\rho^d\) is the chosen extension, the new ray
has primitive lattice vector
\begin{equation}
 v_0=(dx/\tau,1),\qquad
 \Lambda_d=\Z v_0\oplus\ker(e).
 \label{trait:reduced-ray}
\end{equation}
Its last coordinate proves primitivity and the displayed
splitting.  Thus the open ray chart is a torus over
\(\kappa[\rho]\), with \(\rho\) a regular coordinate
vanishing to order one.  Including the free smooth
parameters gives a regular chart with a single reduced
smooth vertical divisor.

Make the algebraic cut-ideal modification over the new
discretely valued function field, and follow the
prolonged \(w\).  Every equal-cut ratio has value zero;
in a blowup chart the ratio or its inverse is regular,
so it is a unit at the centre.  Imposing all these
unit conditions, and removing boundary components
outside the supporting face, puts the centre in this
open ray locus.  This remains true when smooth
parameters vanish additionally at the centre: the
open ray chart, with its smooth parameters, is still
regular.  The vertical prime is the prime with the
label already specified, and its scaled order agrees
with \(w_x\) by \eqref{trait:root-lattice}.  The rescaled
boundary functions \(z'_i\) are units: a suitable
power is a toric character of zero ray order times a
local unit, and normality removes that power.  The
chosen free stratum units remain units as well.

These assertions hold on the algebraic trait, not only
on its geometric completion.  The map from its
excellent characteristic-zero valuation DVR to the
completed trait with algebraically closed residue is
flat and regular, with the same uniformizer.  Regular
base change preserves normality.  Geometric generic
integrality and flatness identify the main component,
so finite main normalization commutes with this base
change.  Blowups commute with flat base change, and
the same normality argument applies to their finite
normalizations.  Smoothness and a reduced smooth
special fibre may consequently be checked in the
geometric completed charts and descend at the centre.
The resulting local vertical divisor downstairs has
the same uniformizer-normalized order.

In this regular local ring factor the principal
divisor of \(H\) into its vertical order and its
horizontal prime factors.  The vertical part gives
\(w_x(H)\).  Thus \(w(H)-w_x(H)\) is a signed sum of
the orders of effective local Cartier horizontal
primes \(D\), with their zero or pole multiplicities.
Their total weighted geometric generic degree is
still at most \(C_1p\).  The modifications have
changed no generic fibre.  Under the possibly
unbounded scalar extension the degrees of split
components sum, with multiplicities, to the original
geometric degree.  Rescaling the \(z_i\) by base
constants changes coefficients, not their rational
presentation degrees.  Thus this degree budget has
no dependence on the auxiliary root degree.

For each horizontal \(D\), choose a nonzero polynomial
\(q(z)\) vanishing on its generic image with
\begin{equation}
 \deg q\le C_3\deg D.
 \label{trait:horizontal-polynomial}
\end{equation}
The \(z\)-projection is defined at the generic point
of \(D\), since the rescaled coordinates are regular
units at our centre and their chosen presentations
have nonvanishing denominators on this neighbourhood.
Its image closure has dimension at most \(n-1\).
Bezout, applied to the bounded-degree presentations
of \(z\), bounds its degree by \(C\deg D\), even
when the image is contracted.  A subvariety of that
degree in affine \(n\)-space lies in a nonzero
hypersurface of at most that degree, by generic linear
projection.  This proves \eqref{trait:horizontal-polynomial}
over the actual extended base field.  Rewrite \(q\)
in \(z'\) and multiply by a base constant so that its
minimum coefficient valuation is zero.  It is regular
at the centre and its reduction is a nonzero polynomial.
It still vanishes along \(D\), whence
\begin{equation}
 0\le w(D)\le v(q).
 \label{trait:horizontal-order}
\end{equation}

On a smooth torus over a characteristic-zero field,
a nonzero polynomial of degree \(d_q>0\) has log
canonical threshold at least \(1/d_q\).  One may
extend the field algebraically and prove this local
assertion at each zero as follows.  Choose successive
general affine linear sections through the point,
ending in a line on which the polynomial is not
identically zero.  Its vanishing order on that line
is at most \(d_q\), so the line pair with coefficient
\(1/d_q\) is lc.  Smooth-divisor inversion of
adjunction \cite{KawakitaInversion}, applied successively
and then dropping the added section divisor, gives
log canonicity in the ambient smooth space near the
point.  Points where the polynomial is a unit impose
no condition.  A nonzero constant likewise contributes
zero order.

Apply this bound to the reduction of the normalized
\(q\).  Inversion of adjunction along the smooth
reduced fibre \(P_0\) gives
\begin{equation}
 v(q)\le(\deg q)A_{T,P_0}(v)
       =(\deg q)E(w).
 \label{trait:polynomial-excess}
\end{equation}
It suffices to apply this over the generic point of
the extracted base prime, where all centres in use
lie.  Spreading its residue-field resolution, or
spreading the polynomial lc assertion over a base
neighbourhood, verifies the hypotheses there.
Combining \eqref{trait:horizontal-polynomial}--\eqref{trait:polynomial-excess}
and summing with the zero and pole multiplicities
proves the first inequality of \eqref{trait:TV}.

For the second inequality keep one logarithmic
coordinate chart and its label along the segment.
There is a polynomial relation
\begin{equation}
 \sum_j a_j(z)H^j=0,\qquad \deg_z a_j\le C_4p.
 \label{trait:relation}
\end{equation}
Indeed, \(z\) is generically finite by differential
independence.  Take the equation of its irreducible
graph image in \(\mathbb P^n_z\times\mathbb P^1_H\),
then dehomogenize.  Its degree in the \(z\)-block
is counted by \(n-1\) general \(z\)-hyperplanes
and one general \(H\)-hyperplane.  These intersections
can be taken in the rational-map open over which
the generic map to the graph image is finite.  The
\(z\) cuts have bounded degree and the general
\(H\) cut has degree at most \(C_1p\), bounded by
its pole divisor.  Bezout in the fixed bounded
embedding proves \eqref{trait:relation}.  The chosen
chart-presentation open is dense on the generic
fibre and its image contains a dense open of the
graph image, so it suffices for this computation.
For constant \(H\) the assertion is immediate.

By the Gauss rule over the fixed \(u\), the value
of each nonzero \(a_j(z)\) is a finite minimum of
affine functions on the segment, with Lipschitz
constant at most \(C_5p\).  In \eqref{trait:relation}
the least term value must be attained at least
twice.  Hence \(w_x(H)\) lies among the finitely
many candidates
\[
 \frac{w_x(a_i)-w_x(a_j)}{j-i}\quad(i\ne j),
\]
each piecewise affine with Lipschitz constant at
most \(2C_5p\).  The actual value is continuous
on the closed labelled segment by the monoid-series
calculation.  Partition the segment at the finitely
many affine breakpoints and intersections of these
candidates.  On each remaining interval a
continuous selection follows one of the affine
candidates, and therefore has the same Lipschitz
bound.  Adding lengths gives that bound on the
whole segment, including its labelled endpoints.
Enlarge \(C_2\) to dominate the constants in the
two inequalities.  This proves \eqref{trait:TV}.
\end{proof}

\subsection{The bounded-index base form}

We prove Theorem~\ref{thm:bounded-fibre}(1).  For this construction
return to the prepared projective families of
Proposition~\ref{bounded:preparation}, before the axis extensions.
The minimum identity \eqref{trait:T} holds on these families by
\eqref{bounded:finite-extension} and the unchanged pulled-back order
functions.  It applies to every allowed field-valued parameter map,
including a nondominant one, and to the universal generic family with
any absolute base volume form, without requiring absolute log canonicity.

Apply the lc-trivial canonical bundle formula and b-semiampleness
\cite{BFMT} to the universal generic form
\(\eta\otimes\zeta^{\otimes r_0}\).  Use a projective fibration model
carrying the original effective lc generic pair; effectivity of its
crepant transform on \(Y\) is unnecessary.  The rank-one condition and
threshold interpretation are precisely those verified in
Section~\ref{sec:tower}.  On every smooth higher base model,
\eqref{trait:T} identifies the discriminant trace as the negative of
the canonical divisor represented by \(\zeta\), minus the pullback of
\(M_0\).  The absolute adjoint is \(\Q\)-principal.  Thus \(M_0\)
represents the actual moduli divisor and is \(\Q\)-semiample.  If the
base-free multiple is first obtained on a higher model, it descends
along the projective birational map: the pushforward of its structure
sheaf is \(\cO_B\), and the pullback divisor is the pullback of
\(M_0\).  Choose \(m\ge2\) with \(mM_0\) globally generated.  The
finitely many packages make this choice uniform.

Resolve the actual parameter map and the divisor supports on a smooth
projective model of \(K\).  The pullback \(M_{0,K}\) is defined even for
a nondominant map, because its generic point lies off the boundary
support.  Formula~\eqref{trait:T} says
\[
 \bar A_\sigma(u)=A_\zeta(u)+u(M_{0,K})\ge0.
\]
Choose a general member
\(H=mM_{0,K}+\operatorname{div}(h)\), transverse to a log resolution of
this lc form boundary.  Along every component of \(H\) the old
boundary coefficient is zero; the reduced transverse addition
therefore preserves log canonicity.  In valuation language,
\(u(H)\le\bar A_\sigma(u)\).  The absolute form
\begin{equation}
 \sigma_K=\frac{\zeta^{\otimes m}}{h}
 \quad\text{satisfies}\quad
 A_{\sigma_K}(u)=\bar A_\sigma(u)-\frac{u(H)}m.
 \label{bounded:base-form}
\end{equation}
Since \(m\ge2\), this lies between
\(\tfrac12\bar A_\sigma(u)\) and \(\bar A_\sigma(u)\).
Taking the norm if a bounded preliminary field extension was used
proves the required bounded index on the original base.  If \(K=\C\),
there are no nonzero base divisorial valuations and a nonzero scalar
form of index one gives the same assertion.

\subsection{Return to bounded denominators}

We prove Theorem~\ref{thm:bounded-fibre}(2).  Let \(u\)
be integral-valued and choose a rational divisorial
\(w\) over it with \(A_\psi(w)\) sufficiently small.
Perform only the bounded parameter and axis
extensions, retaining the old scale.  Their degrees
are bounded by some fixed integer \(N\).  Set
\(D_0=\operatorname{lcm}(1,\ldots,N)\), enlarging
\(N\) to bound the total extension degree.  The
values \(b,\tau\) then belong to
\(D_0^{-1}\Z\).  All data in this paragraph refer
to these bounded extensions; the auxiliary roots
used to prove \eqref{trait:TV} have been discarded.

Equations \eqref{trait:T} and \eqref{trait:TE} give
\[
 A_\sigma(w)=\bar A_\sigma(u)
             +\bigl(\phi(x)-\mu(b)\bigr)+E(w).
\]
The first two summands are nonnegative.  The
hypothesis \(A_\sigma\le C A_\psi\) therefore yields
\(E(w)\le C A_\psi(w)\).  Applying
\eqref{trait:TV} to the ratio with \(\psi\) gives
\begin{equation}
 A_\psi(w_x)\le C_6 A_\psi(w),\qquad C_6=1+CC_2.
 \label{trait:small-retraction}
\end{equation}
In fact the difference equals \(-E(w)\) plus
\((w_x(H)-w(H))/p\); the displayed larger constant
works without a sign restriction on \(C_2-1\).
On each closed labelled cell, \(A_\psi(w_x)\) has
a uniform Lipschitz constant \(L_0\): this follows
from the second inequality of \eqref{trait:TV} and
the fixed linear part \(\phi\) in \eqref{trait:TE}.

Write \(x=\sum_{i=1}^s d_i r_i\), with
\(d_i\ge0\) and the primitive rays \(r_i\) of its
source cone.  For the finite list of cones choose
uniform bounds
\[
 S_0\ge s,\qquad
 P_0'\ge\sum_i\|p_*r_i\|_\infty,\qquad
 R_0\ge\sum_i\|r_i\|.
\]
Given a small \(\epsilon>0\), choose an integer
\(Q\ge2\) with \(Q>P_0'\) and \(Q^{-1}\le\epsilon\).
Place the \(Q^s+1\) points
\[
 j(D_0d_1,\ldots,D_0d_s)\pmod{\Z^s},
 \qquad 0\le j\le Q^s,
\]
in the \(Q^s\) half-open cubes of side \(Q^{-1}\).
Two lie in the same cube.  Subtracting them gives
\(1\le h\le Q^s\) and integers \(q_i\) such that,
with \(m=D_0h\),
\begin{equation}
 |md_i-q_i|<Q^{-1},\qquad
 D_0\le m\le D_0Q^{S_0}.
 \label{trait:dirichlet}
\end{equation}
Since \(md_i\ge0\) and \(Q^{-1}<1\), all the
\(q_i\) are nonnegative.  The residual vector
\(\sum_iq_ip_*r_i-mb\) is integral, and its sup
norm is less than \(P_0'/Q<1\).  Hence it is
exactly zero.  Thus
\[
 y'=\sum_i\frac{q_i}{m}r_i
 \quad\text{satisfies}\quad
 p_*y'=b,\qquad
 m\|y'-x\|<R_0\epsilon.
\]
It lies in the same closed cell, with the incident
label supplied by the fixed chart, even if some
\(q_i\) are zero.  For a cone without rays the
same assertion holds with \(x=y'=0\) and
\(m=D_0\); the approximation step is empty.

Now \((my',m\tau)\in\Lambda_\delta\), and
\(m\tau>0\).  By \eqref{trait:integral-weight}
and the Gauss residue calculation,
\(v=mw_{y'}\) is integral-valued and divisorial.
Its restriction to the original field remains
integral-valued and divisorial, and its base
restriction is \(mu\).  Homogeneity, the
Lipschitz estimate and \eqref{trait:small-retraction}
give
\begin{equation}
 A_\psi(v)\le mC_6 A_\psi(w)+L_0R_0\epsilon.
 \label{trait:final-lift}
\end{equation}
First choose \(\epsilon\) with
\(L_0R_0\epsilon<\lambda/2\), omitting this
condition if \(L_0R_0=0\).  Then choose the input
threshold smaller than
\(\lambda/(2D_0Q^{S_0}C_6)\).  If
\(\bar A_\psi(u)\) is below that threshold,
the defining infimum supplies the required
\(w\), and \eqref{trait:final-lift} gives
\(A_\psi(v)<\lambda\).  The choices are in the
stated order and are uniform.  This proves the
second assertion of Theorem~\ref{thm:bounded-fibre}.
